\subsection{弱空间的商空间与子空间}

设 F、G 为处于对偶的两个实向量空间。设 M 为 F 的一个向量子空间，并考察 G 中正交子空间 \(M^{\circ}\) 内的子空间 N；若 \(y_{1}, y_{2}\) 是 G 中模 N 同余的两点，则 \(\langle x, y_{1} \rangle = \langle x, y_{2} \rangle\) 对所有 \(x \in M\) 成立。对每个模 N 的类 \(\dot{y}\)，把 \(\langle x, y \rangle\) 当 y 取遍 \(\dot{y}\) 时所取的公共值记作 \(\langle x, \dot{y} \rangle\)；显然 \((x, \dot{y}) \mapsto \langle x, \dot{y} \rangle\) 是 \(M \times (G/N)\) 上的一个双线性型。

命题 7. --- 设 M 为 F 的一个向量子空间、N 为 G 的一个向量子空间，其中 F 与 G 是处于对偶的两个向量空间。假设 M 与 N 正交（这等价于 \(N \subset M^{\circ}\)，或 \(M \subset N^{\circ}\)）。则向量空间 M 与 G/N 按双线性型 \((x, \dot{y}) \mapsto \langle x, \dot{y} \rangle\) 处于对偶。

\begin{enumerate}
\def\labelenumi{(\roman{enumi})}
\item
  对这一对偶而言，拓扑 \(\sigma(\mathbf{M},\mathbf{G} / \mathbf{N})\) 由 \(\sigma(\mathbf{F},\mathbf{G})\) 诱导（特别地，我们有 \(\sigma(\mathbf{F},\mathbf{G}) = \sigma(\mathbf{F},\mathbf{G} / \mathbf{F}^{\circ})\)）。
\item
  对这一对偶而言，拓扑 \(\sigma(\mathbf{G}/\mathbf{N}, \mathbf{M})\) 粗于 \(\sigma(\mathbf{G}, \mathbf{F})\) 关于 \(\mathbf{N}\) 的商拓扑；这两个拓扑相同当且仅当 \(\mathbf{M} + \mathbf{G}^{\circ} = \mathbf{N}^{\circ}\)。
\setcounter{enumi}{0}
\item
  \(\mathbf{G} / \mathbf{N}\) 的每个元素都是模 \(\mathbf{N}\) 的一个类，由 \(\mathbf{G}\) 中的某个元素给出；若 \(z_{i}\)（\(1 \leqslant i \leqslant n\)）是 \(\mathbf{G}\) 的元素，且 \(\dot{z}_{i}\)（\(1 \leqslant i \leqslant n\)）是 \(z_{i}\) 在 \(\mathbf{G} / \mathbf{N}\) 中的类，则 \(y \in \mathbf{M}\) 中使 \(|\langle y, \dot{z}_{i} \rangle| \leqslant \alpha\) 对 \(1 \leqslant i \leqslant n\) 成立的 y 所成之集，是 \(\mathbf{M}\) 上这样的 \(x \in \mathbf{F}\)——即满足 \(|\langle x, z_{i} \rangle| \leqslant \alpha\)（\(1 \leqslant i \leqslant n\)）者——所成之集的迹；结论由弱拓扑零邻域的定义即得。
\item
  设 \(p: \mathbf{G} \to \mathbf{G} / \mathbf{N}\) 为典范满射。我们证明，商拓扑 \(\mathcal{T}\)——由 \(\sigma(\mathbf{G}, \mathbf{F})\) 按 \(\mathbf{N}\) 取商而得——与 \(\sigma(\mathbf{G} / \mathbf{N}, \mathbf{N}^{\circ})\) 一致。由于对 \(z \in \mathbf{G}\)、\(y \in \mathbf{N}^{\circ}\) 有 \(\langle y, p(z) \rangle = \langle y, z \rangle\)，可见 \(\sigma(\mathbf{G} / \mathbf{N}, \mathbf{N}^{\circ})\) 的每个零邻域都具有 \(p(\mathbf{V})\) 的形式，其中 \(\mathbf{V}\) 是 \(\sigma(\mathbf{G}, \mathbf{F})\) 的一个对关系 \(z - z' \in \mathbf{N}\) 饱和的零邻域，因此 \(\mathcal{T}\) 细于 \(\sigma(\mathbf{G} / \mathbf{N}, \mathbf{N}^{\circ})\)。 反之，设 \(\mathbf{U} = \mathbf{W}(y_1, ..., y_n; \alpha)\) 为 \(\mathbf{G}\) 中 \(\sigma(\mathbf{G}, \mathbf{F})\) 的一个零邻域，其中 \(y_i \in \mathbf{F}\)（\(1 \leqslant i \leqslant n\)），且 \(\alpha > 0\)；我们将看到，对 \(1 \leqslant i \leqslant n\)，存在元素 \(t_i \in \mathbf{N}^{\circ}\)，使得若令 \(\mathbf{U}' = \mathbf{W}(t_1, ..., t_n; \alpha)\)，则有 \(p(\mathbf{U}') \subset p(\mathbf{U})\)；这就表明 \(\sigma(\mathbf{G} / \mathbf{N}, \mathbf{N}^{\circ})\) 细于 \(\mathcal{T}\)，从而确实与 \(\mathcal{T}\) 一致。 现在，设 \(L\) 为 \(F\) 中由 \(N^{\circ}\) 与诸 \(y_i\) 生成的向量子空间，并用 \(P\) 表示 \(N^{\circ}\) 在 \(L\) 中的补子空间；它是有限维的，设其维数为 \(m\)。设 \((x_j)_{1 \leqslant j \leqslant m}\) 为 \(P\) 的一组基；限制到 \(N\) 上的线性型 \(x \mapsto \langle x_j, z \rangle\) 是线性无关的，因为否则便存在 \(x \neq 0\)（它属于 \(P\)），使得 \(\langle x, z \rangle = 0\) 对一切 \(z \in N\) 成立，也就是说 \(x \in N^{\circ}\)，这与 \(P\) 的定义矛盾。由此我们得出，对所有 \(z' \in G\)，存在 \(s \in N\) 使得 \(\langle x_j, z' \rangle = \langle x_j, s \rangle\) 对所有 \(j\) 成立；若 \(z' = z + s\)，则 \(\langle x, z \rangle = 0\) 对所有 \(x \in P\) 成立。作了这样的准备之后，写 \(y_i = t_i + w_i\)，其中 \(t_i \in N^{\circ}\)、\(w_i \in P\)；我们有 \(\langle y_i, z \rangle = \langle t_i, z \rangle = \langle t_i, z' \rangle\)（\(1 \leqslant i \leqslant n\)）；于是，对所有 \(z' \in U'\)，存在 \(z \in U\) 使得 \(z' - z \in N\)，也就是说有 \(p(\mathrm{U}') \subset p(\mathrm{U})\)。
\end{enumerate}

回到 M 是 \(N^{\circ}\) 中任意子空间的情形，注意显然有 \(\sigma(\mathrm{G}/\mathrm{N},\mathrm{M})=\sigma(\mathrm{G}/\mathrm{N},\mathrm{M}+\mathrm{G}^{\circ})\)；此外，由 II, p.~43 的命题 3 可见，若 \(y\in N^{\circ}\) 使得线性型 \(\dot{z}\mapsto\langle y,\dot{z}\rangle\) 关于 \(\sigma(\mathrm{G}/\mathrm{N},\mathrm{M})\) 连续，则必有 \(y\in M+G^{\circ}\)。我们由此断定：条件 \(M+G^{\circ}=N^{\circ}\) 是商拓扑 T 等于 \(\sigma(\mathrm{G}/\mathrm{N},\mathrm{M})\) 的充分必要条件。

注. --- M 与 G/N 之间的对偶（其中 M 与 N 是两个正交的子空间）在 M 中是分离的，当且仅当 \(M \cap G^{\circ} = \{0\}\)；它在 G/N 中是分离的，当且仅当 \(N = M^{\circ}\)。

推论 1. --- 假设 F 与 G 之间的对偶在 F 中是分离的。对 F 的向量子空间 M 而言，拓扑 \(\sigma(\mathbf{G}/\mathbf{M}^{\circ}, \mathbf{M})\) 与 \(\sigma(\mathbf{G}, \mathbf{F})\) 关于 \(\mathbf{M}^{\circ}\) 的商拓扑一致，当且仅当 M 关于拓扑 \(\sigma(\mathbf{F}, \mathbf{G})\) 是闭的。

这可由命题 7 取 \(\mathbf{N} = \mathbf{M}^{\circ}\)，并回忆 \(\mathbf{M}^{\circ \circ}\) 是 \(\mathbf{M}\) 关于 \(\sigma (\mathrm{F},\mathrm{G})\) 的闭包而得到（II, p.~45, 推论 2）。

推论 2. --- 若 M 是 n 维有限的且对偶在 F 中分离，则 \(M^{\circ}\) 在 G 中的余维数是 n。若 M 关于 \(\sigma(F, G)\) 是闭的且具有有限余维数 n，且对偶在 G 中分离，则 \(M^{\circ}\) 的维数是 n。

因为 \(G/M^{\circ}\) 与 M 处于分离对偶；若 M 的维数是 n，那么 \(G/M^{\circ}\) 亦然（II, p.~41, 例 1）。若 M 是闭的，则 \(F/M = F/M^{\circ\circ}\) 与 \(M^{\circ}\) 处于分离对偶；若 F/M 的维数是 n，那么 \(M^{\circ}\) 的维数也是 n（II, p.~41, 例 1）。

推论 3. --- 设 (F, G)、(F₁, G₁) 是两对处于分离对偶的空间，u 是 F 到 F₁ 中的线性映射，它关于 σ(F, G) 与 σ(F₁, G₁) 连续。则 u 是 F 到 F₁ 中的严格态射，当且仅当 Im(\({}^t u\)) 是 G 中关于 σ(G, F) 的闭子空间。

设 \(\mathbf{N} = \operatorname{Im}(^t u) \subset \mathbf{G}\)；我们知道 \(\mathrm{N}^{\circ} = \mathrm{Ker}(u)\) 在 \(\mathbf{F}\) 中成立（II, p.~47, 公式 (3)）。设 \(p: \mathbf{F} \to \mathbf{F}/\mathbf{N}^{\circ}\) 为典范映射，于是 \(u\) 分解为

\[
u: \mathrm{F} \xrightarrow {p} \mathrm{F} / \mathrm{N} ^ {\circ} \xrightarrow {w} \mathrm{F} _ {1},
\]

其中 w 是单射。空间 \(F/N^{\circ}\) 与 N 处于分离对偶，且由 II, p.~48 的公式 (1)，有 \(\langle w(\dot{y}), z_{1} \rangle = \langle \dot{y}, {}^{t}u(z_{1}) \rangle\)（对所有 \(\dot{y} \in F/N^{\circ}\) 与 \(z_{1} \in G_{1}\)）。这个关系式表明 w 是 \(F/N^{\circ}\)（赋拓扑 \(\sigma(F/N^{\circ}, N)\)）到 \(u(F)\)（赋 \(\sigma(F_{1}, G_{1})\) 诱导的拓扑）上的同构。因此结论由推论 1 以及严格态射的定义得出。

推论 4. --- 设 \((\mathrm{F}, \mathrm{G})\)、\((\mathrm{F}_{1}, \mathrm{G}_{1})\) 是两对处于分离对偶的空间，u 是 F 到 \(F_{1}\) 中的线性映射，它关于 \(\sigma(\mathrm{F}, \mathrm{G})\) 与 \(\sigma(\mathrm{F}_{1}, \mathrm{G}_{1})\) 连续。则 u 是满射，当且仅当 \(^{t}u\) 是 \(G_{1}\)（赋拓扑 \(\sigma(G_{1}, F_{1})\)）到 \(^{t}u(G_{1})\)（赋 \(\sigma(\mathrm{G}, \mathrm{F})\) 诱导的拓扑）上的同构。

因为 \(u(\mathrm{F}) = \mathrm{F}_{1}\) 等价于 \(u(\mathrm{F})\) 在 \(F_{1}\) 中关于 \(\sigma(\mathrm{F}_{1}, \mathrm{G}_{1})\) 既闭又处处稠密；于是推论 4 便由应用于 \(^{t}u\) 的推论 3 以及 II, p.~47 的推论 2 得出。

注. --- 1) 设 \((\mathrm{F}_1, \mathrm{G}_1), (\mathrm{F}_2, \mathrm{G}_2), (\mathrm{F}_3, \mathrm{G}_3)\) 为三对处于分离对偶的空间，考虑由两个线性映射组成的序列

\[
\mathrm{F} _ {1} \xrightarrow {u} \mathrm{F} _ {2} \xrightarrow {v} \mathrm{F} _ {3}\tag{5}
\]

它们分别关于相应于 \(\mathbf{G}_1\)、\(\mathbf{G}_2\)、\(\mathbf{G}_3\) 的弱拓扑连续；我们考虑转置映射的序列

\[
\mathrm{G} _ {3} \stackrel {{t _ {v}}} {{\to}} \mathrm{G} _ {2} \stackrel {{t _ {u}}} {{\to}} \mathrm{G} _ {1}.\tag{6}
\]

显然 \({}^t (v\circ u) = {}^t u\circ {}^tv\)，因此关系式 \(v\circ u = 0\) 等价于 \({}^t u\circ {}^tv = 0\)。序列 (5) 正合当且仅当下述三个条件成立：a) \({}^t u\circ {}^tv = 0\)；

\begin{enumerate}
\def\labelenumi{\alph{enumi})}
\setcounter{enumi}{1}
\item
  \(\operatorname{Im}(^{t}v)\) 在 \(\operatorname{Ker}(^{t}u)\) 中稠密；
\item
  \(^t u\) 是 \(G_2\) 到 \(G_1\) 中的严格态射。
\end{enumerate}

这事实上由 II, p.~49 的推论 3 以及 II, p.~47 的公式 (3) 与 (4) 得出。2) 切不可认为当 \(u\) 是 \(F\) 到 \(F_1\) 中的严格态射时，\(^t u\) 就必然是 \(G_1\) 到 \(G\) 中的严格态射；换言之，\(u\) 可以是严格态射，而 \(u(F)\) 却在 \(F_1\) 中关于 \(\sigma(F_1, G_1)\) 不闭。取 \(F\) 为 \(F_1\) 的一个非闭子空间且 \(G = G_1 / F^\circ\)、\(u\) 为典范单射，这一例子便说明了这一点。类似地，序列 (5) 正合这一事实并不必然蕴含 (6) 正合；然而，若序列 (5) 正合且 \(v\) 是严格态射，则序列 (6) 正合，这由注 1 以及 II, p.~49 的推论 3 得出。

